\section{Механизмът в octet}
\label{sec:mechanism}

Този раздел описва как \code{octet} обслужва една връзка, от приемането до отговора.
Откъсите от кода се четат при компилиране на документа направо от изходните файлове на
хранилището \cite{octet}, затова текстът не може да се разминава с кода. Номерата на
редовете са номерата във файла.

\subsection{Пътят на една връзка}

Сървърът приема връзките в един цикъл и пуска за всяка от тях отделна нишка. След това
обработката минава през четири стъпки \figref{fig:path}. Нишката чете първите налични
октети, разпознава протокола по първия от тях и подава връзката на TLS слоя или направо
на анализатора на HTTP. Всичко, което не е TLS или HTTP, се затваря без отговор.

\begin{figure}[H]
\centering
\begin{tikzpicture}[node distance=0.55cm and 0.6cm]
\node[box] (acc) {приемане\\\code{accept()}};
\node[box, below=of acc] (read) {четене на първите\\октети, до 2048};
\node[box, below=of read] (cls) {разпознаване\\по първия октет};
\node[good, below left=0.8cm and 0.4cm of cls] (tls) {TLS ръкостискане\\(OpenSSL)};
\node[good, below=of cls] (http) {анализатор\\на HTTP};
\node[bad, below right=0.8cm and 0.4cm of cls] (unk) {затваряне\\без отговор};
\node[good, below=of tls] (http2) {анализатор на HTTP\\над TLS};
\node[box, below=2.3cm of http] (resp) {отговор и затваряне};
\draw[arr] (acc) -- (read);
\draw[arr] (read) -- (cls);
\draw[arr] (cls) -- node[left, font=\small, xshift=-2pt] {0x16} (tls);
\draw[arr] (cls) -- node[right, font=\small] {A\dots Z} (http);
\draw[arr] (cls) -- node[right, font=\small, xshift=2pt] {друго} (unk);
\draw[arr] (tls) -- (http2);
\draw[arr] (http2) |- (resp);
\draw[arr] (http) -- (resp);
\draw[thick, gray, |-|] ([xshift=-2.6cm]read.north -| tls.west) -- ([xshift=-2.6cm]tls.south -| tls.west)
  node[midway, left, font=\small, align=right, text=black] {един общ\\краен срок};
\end{tikzpicture}
\caption{Пътят на една връзка в \code{octet}}
\label{fig:path}
\end{figure}

Листинг~\ref{lst:handle} показва функцията, която изпълнява този път. Тя е цялата логика
на демултиплексирането, около петдесет реда.

\lstinputlisting[firstline=133, lastline=183, firstnumber=133, caption={Обработката на
една връзка, \code{src/server.cpp}}, label={lst:handle}]{../src/server.cpp}

\subsection{Разпознаването}

Функцията \code{classify} (листинг~\ref{lst:classify}) получава прочетените октети и
връща едно от три решения: TLS, некриптиран HTTP или непознато. Тя гледа само първия
октет. Тя е \code{constexpr} и \code{noexcept}: не заделя памет, не вдига изключения и
може да се изчисли още при компилиране. Тестовете използват това свойство, за да
проверят част от поведението ѝ още при компилирането на теста.

\lstinputlisting[firstline=36, lastline=52, firstnumber=36, caption={Разпознаването,
\code{include/octet/classify.hpp}}, label={lst:classify}]{../include/octet/classify.hpp}

Празен вход дава „непознато“, а не предположение. Това има значение при връзка, която е
затворена, преди да изпрати нещо.

\subsection{Връщане на прочетените октети}

Разпознаването трябва да прочете поне един октет от сокета. Този октет обаче е и първият
октет на ClientHello или на реда на заявката, и слоят след разпознаването се нуждае от
него. Има два начина да не се изгуби.

Първият е флагът \code{MSG\_PEEK} на \code{recv}: октетите се поглеждат, но остават в
буфера на ядрото и следващото четене ги получава отново. Цената е второ системно
извикване, което чете същите данни още веднъж.

Вторият начин е октетите да се прочетат наистина и после да се подадат на следващия
слой, преди той да чете от сокета. \code{octet} използва втория. Класът
\code{ReplayReader} (листинг~\ref{lst:replay}) държи прочетените октети и при всяко
четене първо връща тях. Едва когато свършат, той чете от сокета. Така никой слой след
разпознаването не може да получи потока разбъркан или с пропуск.

\lstinputlisting[firstline=38, lastline=50, firstnumber=38, caption={Четенето с връщане,
\code{include/octet/replay.hpp}}, label={lst:replay}]{../include/octet/replay.hpp}

При разпознаването сървърът иска до 2048 октета, а не само един. Това не струва нищо
повече, защото системното извикване е едно и също, а обикновено носи целия ClientHello
или целия ред на заявката. При всички \winTotal{} измерени TLS връзки с разпознаване във
второто измерване на Windows първото четене е донесло целия ClientHello от
\winFirstReadMedian{} октета (същото е и в другите две измервания). Това означава, че
при тези условия разпознаването не добавя нито едно системно извикване: OpenSSL получава
ClientHello от буфера, а не от сокета. Условията са важни: клиентът и сървърът са на
една машина и ClientHello пристига в един сегмент. Ако той пристигне на части,
механизмът продължава да работи, както показва тестът в раздел~\ref{sec:verification},
но тогава четенията са повече.

\subsection{OpenSSL върху буфери в паметта}

Обикновено OpenSSL чете и пише сокета сам. Тогава обаче не може да получи октетите,
които разпознаването вече е прочело. Затова \code{octet} свързва OpenSSL с два буфера в
паметта (BIO от вида \code{BIO\_s\_mem} \cite{openssl_biomem}). Сървърът пълни входния
буфер с октети от \code{ReplayReader} и изпразва изходния буфер към сокета.

Празен буфер в паметта по подразбиране означава край на потока, което OpenSSL би приел
за прекъсната връзка. С \code{BIO\_set\_mem\_eof\_return(rbio, -1)} празният буфер
вместо това връща „опитай пак“ \cite{openssl_biomem}. Тогава OpenSSL отговаря с
\code{SSL\_ERROR\_WANT\_READ}, сървърът чете още октети и опитва отново.

Листинг~\ref{lst:handshake} показва цикъла на ръкостискането. Всяка стъпка завършва с
изпращане на всичко, което OpenSSL е произвело. При успех това е следващата група
съобщения, а при провал е предупреждението. Предупреждение, което не е изпратено,
оставя клиента да гадае защо връзката е прекъсната.

\lstinputlisting[firstline=227, lastline=268, firstnumber=227, caption={Цикълът на
ръкостискането, \code{src/tls.cpp}}, label={lst:handshake}]{../src/tls.cpp}

Тази схема има и второ предимство: крайният срок остава в кода на сървъра. Всяко четене
от сокета минава през \code{recv\_some}, което чака само толкова, колкото е останало от
срока. OpenSSL никога не блокира в собствено четене.

\subsection{Един общ краен срок}

Относителният срок, който започва отначало при всяко четене, не защитава от клиент,
който изпраща по един октет точно преди изтичането му. Такъв клиент държи връзката
отворена неограничено. Затова \code{octet} използва абсолютен срок
(листинг~\ref{lst:deadline}): един момент във времето, от който всяко чакане изчислява
колко е останало.

\lstinputlisting[firstline=25, lastline=37, firstnumber=25, caption={Оставащото време до
срока, \code{include/octet/deadline.hpp}}, label={lst:deadline}]{../include/octet/deadline.hpp}

Срокът започва при приемането на връзката и покрива и разпознаването, и цялото
ръкостискане (листинг~\ref{lst:handle}, ред 137). Клиент, който изпраща 0x16 и спира,
е същата заплаха като клиент, който не изпраща нищо. Двата получават един и същ срок.
След завършването на ръкостискането започва втори, отделен срок за заглавната част на
HTTP заявката.

\subsection{Настройките на TLS контекста}

Контекстът на OpenSSL се създава веднъж при стартиране. Таблица~\ref{tab:context}
изброява настройките му и причината за всяка.

\begin{table}[H]
\centering
\caption{Настройки на TLS контекста в \code{octet}}
\label{tab:context}
\begin{tabularx}{\textwidth}{@{}L{5.6cm}Y@{}}
\toprule
Настройка & Причина \\
\midrule
минимална версия TLS 1.2, по избор 1.3 & TLS 1.1 и по-старите не трябва да се договарят
  \cite[раздел 3.1.1]{rfc9325} \\
\code{SSL\_OP\_NO\_RENEGOTIATION} & предоговарянето в TLS 1.2 е било обект на атаки
  \cite{rfc5746}, \cite[раздел 2.10]{rfc7457}; в TLS 1.3 то не съществува \\
\code{SSL\_OP\_NO\_COMPRESSION} & компресията на ниво TLS отваря атаки като CRIME
  \cite[раздел 2.6]{rfc7457}; препоръката е да не се поддържа \cite[раздел 3.3]{rfc9325} \\
\code{SSL\_OP\_CIPHER\_\allowbreak SERVER\_PREFERENCE} & при TLS 1.2 сървърът, а не
  клиентът, избира шифъра от общите \cite{openssl_options} \\
строг ALPN & без общ приложен протокол връзката се отказва с предупреждение 120
  \cite[раздел 3.8]{rfc9325} \\
проверка, че ключът съответства на сертификата & грешна конфигурация спира сървъра при
  стартиране, а не при първия клиент \\
\bottomrule
\end{tabularx}
\end{table}

Изборът на ALPN е в листинг~\ref{lst:alpn}. Ако клиентът предлага \code{http/1.1},
сървърът го избира. Ако не го предлага, строгият режим връща
\code{SSL\_TLSEXT\_ERR\_ALERT\_FATAL}, което според документацията на OpenSSL означава,
че няма общ протокол \cite{openssl_alpn}. В раздел~\ref{sec:verification} се показва, че
OpenSSL тогава наистина изпраща предупреждение 120. Клиент, който изобщо не изпраща ALPN
разширение, се приема, защото няма какво да се откаже.

\lstinputlisting[firstline=22, lastline=37, firstnumber=22, caption={Изборът на
приложен протокол, \code{src/tls.cpp}}, label={lst:alpn}]{../src/tls.cpp}

\subsection{Тестовият удостоверяващ орган}

За тестовете и за измерването \code{octet} създава собствена малка инфраструктура на
публичния ключ. Програмата \code{octet-mkcert} генерира два ключа ECDSA върху кривата
P-256 и два сертификата. Първият е на удостоверяващия орган: основно ограничение
CA:TRUE с дължина на пътя 0 и употреба на ключа само за подписване на сертификати и
списъци за отнемане. Вторият е на сървъра, подписан от органа: CA:FALSE, употреба за
цифров подпис, разширена употреба за удостоверяване на сървър и алтернативни имена
\code{localhost} и \code{127.0.0.1}. Срокът е 30 дни.

Две нива вместо един самоподписан сертификат има по конкретна причина. Ако клиентът
вярва директно на ключа на сървъра, не се проверява верига. Тук клиентът вярва на органа
и сървърът трябва да докаже името си със сертификат, подписан от него. Това е веригата,
която проверява и браузърът. Тестовете използват стандартната проверка на клиента,
включително проверката на името, без да я изключват. Ключовете са само за тестове и не
се пазят в хранилището.
